%% ============================================================
%% CHAPTER 1 — INTRODUCTION
%% ============================================================
\chapter{Introduction}
\label{ch:introduction}

\section{Background and Context}
\label{sec:background}

Somewhere between 20 and 50\,\% of the fuel burned in industrial facilities
never does useful work. It departs the process as hot exhaust gas, cooling
water, or radiated heat and is rejected to the
atmosphere~\cite{Johnson2008,Papapetrou2018}. Across manufacturing and
processing in aggregate --- cement kilns operating above 900\,\textdegree C,
electric arc steel furnaces, gas compression trains, refinery distillation
columns --- these losses are large enough that converting even a modest
fraction into electricity would materially reduce both operating costs and
carbon dioxide emissions without requiring new primary energy
infrastructure~\cite{Bonilla2015}.

Australia has particular reason to take industrial waste heat seriously. Cement
manufacture, liquefied natural gas (LNG) processing, and electric arc furnace
steelmaking rank among the country's most energy-intensive industries, all
operating at temperatures at which power cycle generation is thermodynamically
viable. The Australian Renewable Energy Agency's industrial transformation
programme identified waste heat recovery as one of the most cost-competitive
near-term decarbonisation options available to heavy
industry~\cite{ARENA2023}. The complication is that the majority of published
engineering analysis of WHR power cycles was conducted in the United States
or Europe, with economic boundary conditions --- electricity prices, ambient
temperatures, labour rates, and equipment supply chains --- that do not
translate directly to Australian conditions. This creates a gap, and the
present thesis is designed to partially close it.

Three power cycle technologies compete for this application. The \textit{Steam
Rankine Cycle} (SRC) is the mature baseline. Subcritical SRC plants typically
convert 33--38\,\% of heat input to electricity, with supercritical steam
pushing that toward 45--48\,\% at high turbine inlet
temperatures~\cite{Tian2022}. Below approximately 300\,\textdegree C,
efficiency declines steeply as the saturation pressure falls~\cite{Tchanche2011},
which is the thermodynamic opening for CO$_2$-based cycles.

Both CO$_2$ variants exploit the low critical point of carbon dioxide:
31.1\,\textdegree C and 7.38\,MPa. The \textit{Supercritical CO$_2$ (sCO$_2$)
Brayton cycle} maintains CO$_2$ above both values throughout the entire loop,
eliminating phase change entirely and enabling compact turbomachinery. The
cost is that compressing a gas always requires more work than pumping a
liquid; at lower source temperatures, the compressor consumes a prohibitive
share of turbine output. The \textit{Transcritical CO$_2$ (tCO$_2$) Rankine
cycle} handles this by condensing CO$_2$ to a liquid on the low-pressure side
and pumping it above the critical pressure, reducing back-work ratio
significantly. The operational constraint is that condensation requires a heat sink below 31.1\,\textdegree C, which is not always achievable in warmer Australian locations~\cite{Jiang2023,Lee2023}.

This thesis compares all three cycles simultaneously under consistent
Australian-specific assumptions, a combination not previously addressed in the
published literature. The outcome is a practical cycle selection framework for
Australian industrial operators considering WHR investment.

\section{Definitions}
\label{sec:definitions}

The following terms are defined here for unambiguous use throughout.

\begin{description}[leftmargin=2em,style=nextline]\sloppy

\item[\textbf{Rankine Cycle (RC):}]
A closed thermodynamic cycle in which the working fluid is pumped as a
liquid (1$\to$2), heated and evaporated at constant pressure (2$\to$3),
expanded through a turbine (3$\to$4), and condensed back to liquid
(4$\to$1). The SRC uses water operating subcritically~\cite{Moran2018}.

\item[\textbf{Supercritical CO$_2$ (sCO$_2$) Brayton Cycle:}]
A closed power cycle in which CO$_2$ remains above its critical pressure
(7.38\,MPa) and temperature (31.1\,\textdegree C) throughout. No phase
change occurs. Near-critical CO$_2$ is substantially less compressible than
an ideal gas, reducing compressor work relative to a conventional gas
turbine, but back-work ratio remains higher than for liquid-pump
cycles~\cite{Dostal2004,Crespi2017}.

\item[\textbf{Transcritical CO$_2$ (tCO$_2$) Rankine Cycle:}]
A cycle that crosses the CO$_2$ saturation boundary. CO$_2$ condenses to a
subcritical liquid on the low-pressure side, is pumped above critical
pressure, heated supercritically, and expanded. Liquid-phase compression
keeps pump work low (back-work ratio 12--23\,\% across the modelled range).
The requirement for a heat sink below 31.1\,\textdegree C is the main
operational constraint~\cite{Jiang2023,Lee2023}.

\item[\textbf{Pseudo-Critical Temperature ($T_{pc}$):}]
The temperature at which the isobaric specific heat $c_p$ reaches its local
maximum at a given supercritical pressure. CO$_2$ undergoes no phase
transition above $P_c$, but $c_p$ peaks sharply in a narrow band near the
saturation line extension, creating anomalously high heat capacity. At
12.8\,MPa, $T_{pc} \approx 57$\,\textdegree C with peak $c_p =
4.4$\,kJ/(kg$\cdot$K); at 20\,MPa, $T_{pc} \approx 76$\,\textdegree C
with peak $c_p = 2.6$\,kJ/(kg$\cdot$K). This non-uniform thermophysical
behaviour is the central transitional effect addressed in
Sections~\ref{sec:results_transitional} and~\ref{sec:rq2_hx_sizing}.

\item[\textbf{Waste Heat Recovery (WHR):}]
Capturing thermal energy from industrial exhaust or process cooling streams
that would otherwise be rejected to the environment, and converting it to
electricity or useful work~\cite{Papapetrou2018}.

\item[\textbf{Thermal Efficiency ($\eta_{th}$):}]
Net work output divided by total heat input: $\eta_{th} =
\dot{W}_{net}/\dot{Q}_{in}$. A first-law metric that quantifies energy
conversion but reveals nothing about the thermodynamic quality of the
losses.

\item[\textbf{Exergy Efficiency ($\eta_{ex}$):}]
Net work output divided by the exergy of the heat input: $\eta_{ex} =
\dot{W}_{net}/\dot{E}_{x,in}$. A second-law metric that accounts for the
thermodynamic quality of the energy supplied, not merely its magnitude.

\item[\textbf{Back-Work Ratio (BWR):}]
The fraction of gross turbine output consumed by the compression device:
$\mathrm{BWR} = \dot{W}_{pump}/\dot{W}_{turbine}$. For SRC: below 1.5\,\%.
For sCO$_2$ Brayton: 27--80\,\% across the modelled range. For tCO$_2$
Rankine: 12--23\,\% (higher than the literature-reported 2--5\,\% owing to
the $P_{low} = 5$\,MPa baseline; sensitivity analysed in
Section~\ref{sec:plow_sensitivity}).

\item[\textbf{Pinch Point Temperature Difference (PPTD):}]
The minimum local temperature difference between hot and cold streams at
any cross-section within a heat exchanger. PPTD governs the trade-off
between heat recovery and exchanger surface area~\cite{Moran2018}.

\item[\textbf{Levelised Cost of Electricity (LCOE):}]
A lifecycle-averaged cost in AUD/MWh representing the electricity price
at which a project breaks even over its operational life, after accounting
for capital cost, annual operating expenditure, capacity factor, and the
time value of money~\cite{Short1995}.

\item[\textbf{Specific Flow Exergy ($ex$):}]
The maximum work extractable from a fluid stream brought reversibly to
the dead state ($T_0 = 25$\,\textdegree C, $P_0 = 101.325$\,kPa):
$ex = (h - h_0) - T_0(s - s_0)$.

\item[\textbf{Size Parameter (SP):}]
A turbomachinery compactness metric: $\mathrm{SP} =
\dot{V}_{out}^{0.5}/\Delta h_{is}^{0.25}$. A smaller SP indicates a more
compact turbine for a given net power output~\cite{Macchi2017}.

\item[\textbf{Stodola Ellipse Law:}]
An empirical approximation describing how turbine mass flow and pressure
ratio evolve under off-design operation. Here applied in the simplified
form $\eta_{od} = \eta_0 \cdot [1 - 0.3(1-f)^2]$ where $f$ is the
mass-flow fraction~\cite{Marchionni2020}.

\end{description}

\section{Knowledge Gaps}
\label{sec:knowledge_gap}

Five specific gaps in the published literature motivated the design of this
study:

\bigskip
\noindent\textbf{Gap 1 --- No three-cycle benchmark under consistent assumptions.}
Existing studies compare at most two technologies simultaneously, under
different boundary conditions and cost assumptions, making cross-study
inference unreliable~\cite{Jiang2023,Kizilkan2020,Baglietto2023}.

\bigskip
\noindent\textbf{Gap 2 --- Combined exergy and pinch analysis is absent for tCO$_2$.}
The application of combined exergy-pinch methodology to the tCO$_2$ Rankine
cycle in a WHR context has not been reported in the accessible
literature~\cite{Sarkar2009,Nami2017}.

\bigskip
\noindent\textbf{Gap 3 --- Australian industrial conditions are unrepresented
in the techno-economic literature.}
All published LCOE and NPV analyses for CO$_2$ WHR cycles use US or European
economic parameters. Australian electricity prices, ambient temperatures, and
supply chain structures differ sufficiently to make direct transplanting of
overseas conclusions unreliable~\cite{Papapetrou2018,ARENA2023}.

\bigskip
\noindent\textbf{Gap 4 --- Transitional thermophysical effects in tCO$_2$ heat
exchangers are not modelled.}
The sharp $c_p$ peak near the pseudo-critical temperature creates a strongly
non-uniform temperature-difference profile along the heater. Two-state-point
energy balances and two-point LMTD sizing cannot represent this profile, so
the location of the pinch and the required heat exchanger area are
misestimated. This gap was specifically identified by the project supervisor.

\bigskip
\noindent\textbf{Gap 5 --- Off-design data for tCO$_2$ cycles are sparse.}
Characterisation of part-load performance and operational flexibility for
tCO$_2$ Rankine cycles remains limited, with Yang et al.~\cite{Yang2024}
representing one of few recent contributions.

\section{Research Questions and Aim}
\label{sec:rqs}

\begin{mdframed}[linecolor=uqpurple,linewidth=1.5pt,
  backgroundcolor=uqpurple!5,roundcorner=4pt,
  innerleftmargin=12pt,innerrightmargin=12pt,
  innertopmargin=10pt,innerbottommargin=10pt]
\textit{To compare the Supercritical CO$_2$ Brayton cycle and Transcritical
CO$_2$ Rankine cycle against the Steam Rankine Cycle for industrial waste heat
recovery, using thermodynamic modelling and techno-economic analysis, and to
determine under what conditions CO$_2$ cycles represent the better choice in
terms of efficiency, capital cost, and economic return within the Australian
energy transition context.}
\end{mdframed}

\bigskip
\begin{description}[leftmargin=1.5em,style=nextline]
\item[\textbf{RQ1 --- Thermodynamic Performance:}]
How do thermal efficiency, exergy efficiency, and back-work ratio compare
for SRC, sCO$_2$ Brayton, and tCO$_2$ Rankine across WHR source temperatures
of 150--600\,\textdegree C, including sensitivity to recuperator effectiveness
and condensing pressure?

\item[\textbf{RQ2 --- System Compactness and Capital Cost:}]
How compact are the three cycles, and what does this imply for capital cost
(AUD/kW) across the 100\,kW to 5\,MW range at the three Australian case
study conditions?

\item[\textbf{RQ3 --- Australian Techno-Economic Analysis:}]
For three representative Australian industrial sites, what are the LCOE, NPV,
payback period, and probability of economic viability for each cycle under
realistic uncertainty?

\item[\textbf{RQ4 --- Off-Design and Part-Load Performance:}]
When the heat source fluctuates, how does each cycle's efficiency degrade,
and what is the relative operational flexibility of the three technologies?
\end{description}

\section{Scope and Limitations}
\label{sec:scope}

The thermodynamic models are steady-state, single-fluid, zero-pressure-drop
representations of idealised cycles. Real-fluid properties are evaluated using
CoolProp~6.x~\cite{Bell2014} with the HEOS backend, which provides high
accuracy near the CO$_2$ critical point. The WHR heat source is modelled as a
dry ideal gas with constant specific heat; real flue gas composition effects
are not considered. All three cycles use fixed isentropic efficiencies
($\eta_{turbine} = 0.85$, $\eta_{pump/compressor} = 0.80$), a known
simplification explicitly quantified in Section~\ref{sec:discussion_limitations}. Capital
cost correlations from Weiland et al.~\cite{Weiland2019} and Turton et
al.~\cite{Turton2012} carry inherent uncertainty of order $\pm$30\,\% and are
applied below their validated size range; this is explicitly propagated
through Monte Carlo analysis. Off-design modelling
uses the simplified Stodola parabolic efficiency correction rather than
full turbomachinery characteristic maps.

\section{Thesis Structure}
\label{sec:structure}

Chapter~\ref{ch:litreview} reviews the literature on WHR power cycles,
transitional CO$_2$ thermophysics, and Australian industrial context.
Chapter~\ref{ch:method} describes the methodology for all four research
questions. Chapters~\ref{ch:results_rq1}--\ref{ch:results_rq4} report
results for RQ1--RQ4 respectively. Chapter~\ref{ch:discussion} synthesises
across research questions, situates the findings within the broader literature,
and explicitly addresses limitations. Chapter~\ref{ch:conclusions} states the
answers to each research question and identifies priority directions for future
work.

%% ============================================================
%% CHAPTER 2 — LITERATURE REVIEW
%% ============================================================
\chapter{Literature Review}
\label{ch:litreview}

\section{Industrial Waste Heat Recovery}
\label{sec:lit_whr}

The quantity of recoverable industrial waste heat is large. Papapetrou et
al.~\cite{Papapetrou2018} estimated the technically available resource across
European manufacturing at 300--900\,TWh per year, with roughly one-third
arising above 300\,\textdegree C. Johnson et al.~\cite{Johnson2008} placed US
industrial waste heat losses at 20--50\,\% of primary energy input, consistent
with the European estimate, with high-temperature sources arising primarily
from metallurgy and cement, medium-grade from refinery and glass, and
low-grade from compressed gas systems and cooling circuits.

For Australia, the Australian Renewable Energy Agency's (ARENA) industrial
transformation study identified WHR as one of the most cost-competitive
near-term options for decarbonising energy-intensive industry~\cite{ARENA2023}.
The Australian Alliance for Energy Productivity (A2EP)~\cite{A2EP2022}
compiled an industry-specific database confirming WHR opportunities at
Australian cement, LNG, and steel facilities at temperatures closely matching
the three case studies defined for this thesis. A distinguishing feature of
the Australian context is the co-existence of high industrial electricity
prices (particularly in the isolated Pilbara grid at AUD\,130/MWh or
above~\cite{AEMO2024}) with climatic conditions that constrain condenser
temperatures --- a combination that creates simultaneously attractive economics
and operational challenges for low-temperature cycle condensation.

\section{Steam Rankine Cycle}
\label{sec:lit_src}

The SRC has driven industrial power generation since the nineteenth century
and its design practice is codified across a wide capacity range. Tian et
al.~\cite{Tian2022} surveyed SRC across WHR temperatures, finding thermal
efficiencies of 33--38\,\% for subcritical configurations and approaching
48\,\% for supercritical steam above 550\,\textdegree C. The well-documented
limitation below 300\,\textdegree C~\cite{Tchanche2011} --- where falling
saturation pressure constrains maximum turbine inlet conditions --- is the
primary motivation for CO$_2$ cycle investigation in the medium-grade
WHR range.

From a second-law perspective, Moran and Shapiro~\cite{Moran2018} identify
the steam boiler as the dominant SRC loss mechanism. Water boils at a fixed
saturation temperature while the heat source is at a higher temperature, so
a sustained temperature difference is maintained throughout the two-phase
region, generating entropy continuously. CO$_2$ cycles address this by
heating supercritically and non-isothermally, tracking the source temperature
profile more closely. This thermodynamic advantage is the foundation of the
CO$_2$ case, though it must be weighed against differences in component cost
and operational behaviour.

Capital cost is the SRC's competitive strength. Turton et
al.~\cite{Turton2012} document established shell-and-tube boiler cost
correlations that benefit significantly from economies of scale. Baglietto et
al.~\cite{Baglietto2023} confirmed that at larger WHR scales, steam's mature
supply chain erodes the sCO$_2$ capital cost advantage, a finding corroborated
by the present study's RQ2 analysis.

\section{Supercritical CO$_2$ Brayton Cycle}
\label{sec:lit_sco2}

Dostal's 2004 MIT dissertation~\cite{Dostal2004} established the foundational
analysis of near-critical compression and demonstrated the compactness
advantage over steam through detailed cycle modelling. Crespi et
al.~\cite{Crespi2017} reviewed major sCO$_2$ configurations --- simple
recuperated, recompression, split-flow --- and documented thermal efficiencies
of 40--50\,\% at turbine inlet temperatures above 500\,\textdegree C. Ahn et
al.~\cite{Ahn2015} assessed technology readiness, identifying turbomachinery
as the critical engineering challenge and noting the absence of demonstrated
commercial-scale plant.

The WHR case for sCO$_2$ is best illustrated by Kizilkan~\cite{Kizilkan2020},
who found sCO$_2$ producing 9,363\,kW net versus 8,275\,kW for SRC at a
350\,\textdegree C cement WHR source --- a 13.1\,\% improvement at respective
thermal efficiencies of 27.6\,\% and 24.2\,\%. Baglietto et
al.~\cite{Baglietto2023} added the critical qualifier that the sCO$_2$ capital
cost advantage over steam erodes at larger scales due to the immaturity of
the PCHE supply chain.

Below approximately 350\,\textdegree C, sCO$_2$ compressor back-work rises
to 20--40\,\% and continues upward as temperature falls~\cite{Crespi2017}.
This is the primary thermodynamic motivation for the tCO$_2$ cycle.

\section{Transcritical CO$_2$ Rankine Cycle}
\label{sec:lit_tco2}

Replacing gas compression with liquid pumping reduces BWR to the 2--5\,\%
range reported in the literature at moderate-to-high temperatures, while
retaining the supercritical heat addition advantage over SRC. Jiang et
al.~\cite{Jiang2023} reported up to 4.98\,\% better thermal efficiency for a
recompression tCO$_2$ configuration over equivalent sCO$_2$ Brayton, which
informed the working hypothesis of this study. Lee and
Kim~\cite{Lee2023} provided experimental validation of a small-scale tCO$_2$
Rankine system, supporting the modelling assumptions of comparable simulation
work.

Yang et al.~\cite{Yang2024} examined off-design behaviour of tCO$_2$ cycles
for low-grade WHR, quantifying efficiency degradation at part-load that had
not previously been characterised --- directly relevant to RQ4 of this study.
Liu et al.~\cite{Liu2024} demonstrated that optimal high-side pressure
increases with turbine inlet temperature, a result incorporated into the
parametric sweep methodology described in Chapter~\ref{ch:method}. Wang et
al.~\cite{Wang2025} extended this to variable heat source conditions.

A practical constraint specific to Australian conditions deserves emphasis.
The tCO$_2$ cycle requires condensation below 31.1\,\textdegree C. In the
Pilbara in summer (ambient 35--42\,\textdegree C~\cite{BOM2024}) or coastal
Queensland, this requires supplementary mechanical or evaporative cooling,
adding capital cost and parasitic load not accounted for in European or
North American studies. This study treats this constraint explicitly in the
RQ3 economic analysis.

\section{Transitional Thermophysical Effects in CO$_2$ Heat Exchangers}
\label{sec:lit_transitional}

The thermodynamic consequences of the pseudo-critical $c_p$ peak in heat
exchanger design have attracted increasing attention since Dostal's original
analysis~\cite{Dostal2004}. Sarkar~\cite{Sarkar2009} identified the primary
heat exchanger as the dominant exergy destruction site in the sCO$_2$
recompression cycle, attributing this to the strongly non-uniform temperature
difference between source and working fluid along the exchanger length. Nami
et al.~\cite{Nami2017} demonstrated that combined exergy-economic optimisation
surfaces improvement opportunities invisible to first-law analysis.

For the tCO$_2$ cycle specifically, CO$_2$ enters the supercritical heater at
pump exit temperature (24--32\,\textdegree C at 12--20\,MPa), passes through
the pseudo-critical region with its pronounced $c_p$ peak, and emerges at
turbine inlet temperature. Because $c_p$ is largest near $T_{pc}$, a large
fraction of the total heat transfer is absorbed within a narrow temperature
band: the CO$_2$ temperature rises slowly through the pseudo-critical region
and then quickly beyond it. Against a waste-gas stream whose temperature falls
linearly with heat duty, this produces a temperature-difference profile that
is pinched at the terminals and widest in the middle of the exchanger --- the
opposite of the uniform driving force that a two-point LMTD calculation
assumes. The quantitative consequence for heat exchanger sizing, addressed in
Section~\ref{sec:rq2_hx_sizing}, had not been reported in the literature
accessed for this study, representing Gap~4 identified in
Section~\ref{sec:knowledge_gap}.

\section{Capital Cost Correlations for WHR Power Systems}
\label{sec:lit_cost}

The published cost literature for small-scale CO$_2$ power systems is limited
and predominantly US-based. The most complete public set of correlations is
that of Weiland, Lance and Pidaparti~\cite{Weiland2019}, derived from
US~DOE project vendor quotes (2017 USD). It gives the PCHE/recuperator cost as
$C_{HX} = 49.45 \cdot UA^{0.7544} \cdot f_T$, where $UA$ is in W/K and $f_T$
is a temperature penalty factor that exceeds unity above 550\,\textdegree C;
the radial sCO$_2$ turbine cost as $406{,}200 \cdot \dot{W}_{sh}^{0.8}$
($\dot{W}_{sh}$ in MW); and direct air coolers as $32.88 \cdot UA^{0.75}$.
Weiland et al. report uncertainty bands of $-31$\,\% to $+38$\,\% on the
recuperator correlation and $-32$\,\% to $+51$\,\% on the radial turbine
correlation, and the turbine data span 8--35\,MW, so application at the
sub-5\,MW scale of industrial WHR is an extrapolation. Weiland and
Thimsen~\cite{Weiland2016} set out the steady-state modelling assumptions and
constraints (compressor inlet conditions, recuperator limits) that underpin
such cycle models. No comparable public correlation exists for small steam
turbine-generators; the US~DOE industrial benchmark~\cite{DOE2012} of roughly
USD\,900/kW at 150\,kW falling to USD\,200/kW at 2\,MW (equipment only) is
the most widely cited small-scale reference.

Turton et al.~\cite{Turton2012} provide cost functions for conventional
equipment (centrifugal pumps and compressors, shell-and-tube heat exchangers)
in 2001 USD (CEPCI 397), which can be escalated via the Chemical Engineering
Plant Cost Index (CEPCI)~\cite{CEPCI2026}. The 2017 annual CEPCI of 567.5 and
the estimated 2026 value of approximately 830 imply an escalation factor of
$1.46\times$ for the Weiland correlations. Conversion of USD to AUD at the
April 2026 spot rate (1 USD = 1.395 AUD, RBA~\cite{RBA2026}) completes the
localisation.

\section{Off-Design Analysis Methods}
\label{sec:lit_offdesign}

The Stodola ellipse law provides a widely-used approximation for turbine
off-design mass flow as a function of inlet/outlet pressures~\cite{Marchionni2020}.
In its simplified form applicable to efficiency estimation, it approximates
part-load isentropic efficiency degradation as a parabolic function of
mass-flow fraction, consistent with the incidence-loss behaviour of radial
turbines. Marchionni et al.~\cite{Marchionni2020} reviewed off-design
performance of advanced thermodynamic cycles and proposed a multi-attribute
flexibility scoring framework that weights turndown ratio, part-load efficiency
retention, and startup complexity --- the framework adopted for RQ4 in this
study.

Yang et al.~\cite{Yang2024} examined off-design performance of tCO$_2$
cycles specifically and found efficiency degradation at part-load had not
previously been quantified. Their results inform the design-point calibration
of the Stodola correction used in Chapter~\ref{ch:results_rq4}.

\section{Summary and Gap Identification}
\label{sec:lit_summary}

The literature reviewed here is extensive but incomplete in five specific
respects identified in Section~\ref{sec:knowledge_gap}. The thermodynamic
promise of CO$_2$ WHR cycles is well established. What is missing is a single,
unified comparison of all three cycle architectures under consistent Australian-
specific assumptions, incorporating transitional thermophysical effects, full
capital cost estimation, and off-design analysis. The methodology described in
Chapter~\ref{ch:method} addresses all five gaps within a single undergraduate
thesis scope.

